\documentclass[../ICIT-Paper-2024.tex]{subfiles}
\usepackage{listings}
\usepackage{color}

\definecolor{mygreen}{rgb}{0,0.6,0}
\definecolor{mygray}{rgb}{0.5,0.5,0.5}
\definecolor{mymauve}{rgb}{0.58,0,0.82}
\lstset{ %
  backgroundcolor=\color{white},   % choose the background color
  basicstyle=\ttfamily\footnotesize, % size of fonts used for the code
  breaklines=true,                 % automatic line breaking only at whitespace
  captionpos=b,                    % sets the caption-position to bottom
  commentstyle=\color{mygreen},    % comment style
  escapeinside={\%*}{*)},          % if you want to add LaTeX within your code
  keywordstyle=\color{blue},       % keyword style
  stringstyle=\color{mymauve},     % string literal style
  columns=fullflexible,            % ensure that all characters have the same width
}
\begin{document}
\paragraph{}
Genetic algorithm (GA) is a well-known algorithm, which is inspired by the biological evolution process~\cite{kingma2014adam}. Genetic algorithm (GA) is an optimization algorithm that is inspired by natural selection. It is a population-based search algorithm, which utilizes the concept of survival of fittest~\cite{zbMATH03584897} It is used to solve complex problems by iteratively improving a population of candidate solutions. 

Genetic algorithms have been successfully applied to various optimization problems, including parameter tuning, scheduling, routing, and machine learning~\cite{katoch2021review}. They are particularly useful for problems with large search spaces, discontinuous or noisy objective functions, and multiple local optima (1) Some key advantages of genetic algorithms include their ability to explore large search spaces, find globally optimal or near-optimal solutions, and handle complex problems that are not well suited for traditional optimization algorithms~\cite{katoch2021review}. The pseudocode~\ref{al3} shows in detail our algorithm.

\begin{algorithm}[H]
\caption{Genetic Algorithm}\label{al3}
\begin{algorithmic}[1]
\Require Population size $N$, mutation rate $p_m$, crossover type $p_c$, maximum number of generation $MAX$
\State Initialize population of size $N$, iteration counter $t$ = 0
\For{$i$ in range(num\_iterations)}
    \State Evaluate the fitness of each individual in the population
    \State Select the best individual for reproduction based on fitness 
    \While{new population not filled}
        \State Select random pairs of parents in the best individual
        \State Perform crossover between pairs of parents crossover type $p_c$
        \State Perform mutation on offspring with probability $p_m$
        \State Add offspring to new population
    \EndWhile
    \State Increment the current iteration t by 1
    \State Replace the current population with the new population
\EndFor
\State \textbf{return} Best individual found
\end{algorithmic}
\end{algorithm}

The fitness function used in the genetic algorithm is the accuracy evaluation function of predictions. This function is similarly used to evaluate the accuracy of other algorithms.


\begin{lstlisting}[language=Python]
def CUSTOM_ACCURACY (first_y, second_y, y_pred):
    condition = numpy.logical_or(
                    y_pred == second_y, 
                    y_pred == first_y
                    )
    count = sum(condition)
    accuracy = count / len(y_pred)
    return accuracy 
\end{lstlisting}

\end{document}